\documentclass[book.tex]{subfiles}

\begin{document}

\chapter{Complex Analysis}
\label{chap:complex_analysis}
\noindent\rule{11cm}{0.4pt} \\
\textbf{Prerequisites:} \autoref{chap:calculus}, \autoref{chap:distributions} \\
\textbf{Difficulty Level:} ** \\
\noindent\rule{11cm}{0.4pt}
\vspace{5mm}

In this chapter, we will discuss the basics of complex analysis. Tools and tricks from complex analysis will arise often later in this book, in our study of quantum field theory, where it will be necessary to evaluate complex integrals over either Euclidean space $\mathbb{R}^3$, $\mathbb{R}^4$ or over Minkowski space. Hello


\section{Contour Integrals}

Contour integrals provide a methodology for evaluating integrals along curves in the complex plane. We define a curve in the complex plane as a function
\begin{align}
z: [a, b] \to \mathbb{C}
\end{align}
A contour $\gamma$ is defined by a finite set of curves connected pointwise. Let $f: \mathbb{C} \to \mathbb{C}$ be a continuous function. The contour integral is defined as
\begin{align}
\int_\gamma f(z) dz &= \int_a^b f(\gamma(t)) \gamma'(t) dt
\end{align}

%\textcolor{red}{TODO:  Show how to evaluate definite integrals with them}

\section{Cauchy Principal Value}

The Cauchy principal value $\mathcal{P}(\frac{1}{x})$ is a distribution defined as follows. For $u \in C^{\infty}_c(\mathbb{R})$

\begin{align}
    \left [ \mathcal{P} \left ( \frac{1}{x} \right )\right ](u) &= \lim_{\varepsilon \to 0^+} \int_{\mathbb{R} \setminus [-\varepsilon, \varepsilon]} \frac{u(x)}{x} dx  \\
    &= \int_0^{\infty} \frac{u(x) - u(-x)}{x} dx 
\end{align}

\section{Sokhotski-Plemelj Theorem}

Let $f$ be a complex-valued function. Let $a, b \in \mathbb{R}$ be such that $a < 0 < b$. Then

\begin{align}
    \lim_{\varepsilon \to 0^+} \int_{a}^b \frac{f(x)}{x \pm i \varepsilon} = \mp i \pi f(0) + \mathcal{P} \int_a^b \frac{f(x)}{x} dx
\end{align}

Note that this equation is a distributional equation; the quantities above may not be well defined as functions. When $f$ is the Dirac-delta, we obtain the following version of the equation above. 

\begin{align}
    \lim_{\varepsilon \to 0^+} \frac{1}{x \pm i \epsilon} = \mp i \pi \delta(x) + \mathcal{P} \left ( \frac{1}{x} \right )
\end{align}

%TODO: Perhaps this deserves a derivation?

%\section{Exercises}
%\begin{enumerate}
%    \item TODO
%\end{enumerate}

\end{document}